\twocolumn[{%
  \begin{@twocolumnfalse}\maketitle
    {\centering\bfseries Objective\par}
    \begin{enumerate}[label=(\Roman*),ref={[\Roman*]}]
        \item \label[objective]{obj:I} Estimate the probability that in the year \(2023\) the maximum discharge will be greater than the value
           \(C\,[\si{\metre\cubed\per\second}]\), where \(C\) is the centre between \(Q_3\) and the next ordered 
           observed value after the 3\textsuperscript{rd} quartile. That is \[
               C = \frac{X_{[3n/4]} + X_{[3n/4] + 1}}{2}
           \] 
       \item \label[objective]{obj:II} Estimate the probability that in the year \(2023\), the maximum discharge will be greater than 
            \(D\,[\si{\metre\cubed\per\second}]\), where \(D\) is the last ordered observed value 
            \(+\qty{100}{\metre\cubed\per\second}\)
        \item \label[objective]{obj:III} What is the probability that in the year \(2023\) the maximum discharge will be lower than the value 
            \(E\,[\si{\metre\cubed\per\second}]\), where \(E\) is the centre between the 4\textsuperscript{th} and
            5\textsuperscript{th} ordered observation? (The probability that there will be this level of draught)
        \item \label[objective]{obj:IV} What is the minimum threshold \(h\,[\si{\metre\cubed\per\second}]\) such, that the probability that the
            maximum annual discharge will be higher than this threshold in \(2023\) is lower than \(0.1\), \(0.04\),
            \(0.02\), \(0.01\) and \(0.004\)? Find the discharges \(h_{0.1}\), \(h_{0.04}\), \(h_{0.02}\), \(h_{0.01}\)
            and \(h_{0.004}\) for a \(10\)-year, \(25\)-year, \(50\)-year, \(100\)-year and \(250\)-year flood.
        \item \label[objective]{obj:V} What is the probability that the maximum discharge in the year \(2023\) will be greater than any of the
            discharges previously observed? That is, there will be a record flood in \(2023\). How many records in the
            maximum discharge do you anticipate in the upcoming span of \(50\) years?
        \item \label[objective]{obj:VI} Assume that there is a city downstream the river, that disposes of a anti-flood barier able to retain the
            maximum discharge \(C\,[\si{\metre\cubed\per\second}]\) from \Cref{obj:I}. If there was to be a flood
            exceeding this barier \(C\), what is the probability, that the barier will be exceeded by more than 
            \(Z\,[\si{\metre\cubed\per\second}]\), where \(Z\) is \(1/2\) from the difference between the last ordered
            observation and \(Q_3\)? That is \[
                Z = \frac{X_{[n]} - X_{[3n/4]}}{2}
            \] 
   \end{enumerate}
   \hrulefill
\bigskip\nopagebreak
\end{@twocolumnfalse}}]
\section{Introduction}%
\label{sec:introduction}
This is an attempt at predicting the flood hazard of the \emph{Niger River}.

\begin{figure}[H]
    \centering
    \incfig{Map_of_River_Niger}
    \caption{Map of the Niger River from its origin in the Guinea Highlands to the Niger Delta}
    \label{fig:map_of_the_river}
\end{figure}

\subsection{Niger River}%
\label{sub:niger_river}
Niger is the main river of West Africa, with a length exceeding \qty{4100}{km}. Its drainage basin (area of land from
which all of the water ends up in the river) is over \qty{2.1e6}{km^2}. It originates in the Guinea Highlands and flows 
through Mali, Niger and Nigeria, before eventually reaching and discharging at the monumental Niger Delta. This 
trajectory makes it the 3\textsuperscript{rd} longest river in Africa only surpassed by the Congo River and, of course, 
the Nile.

Water flowing through the Niger river is significantly clearer compared to the Nile since it originates in ancient 
rocks providing little clay unlike the Nile. Similarly to Nile however, the Niger River floods annually usually in the 
time period from September to May, peaking in November.

\subsection{Data}%
\label{sub:data}
The data is downloaded from the publicly available \href{https://www.bafg.de/GRDC/EN/Home/homepage_node.html}{\ac{GRDC}}
database. The station, where the measurement took place is located in \emph{Niamey, Niger} (the capital city) at the river 
kilometre \qty{1420}{} (measured from the \emph{headwater} i.e. the furthest point upstream). The station's altitude is \qty{176}{m} above sea level and it covers a basin size of
\qty{791121}{km^2} (i.e. water flowing through this point of the river originates in this amount of area
upstream).

All measurements adhere to the standard notation in the branch, meaning that the discharges are denoted \(Q\). The data
consists of monthly mean measurements and daily measurements.

Monitoring at this station started in the year 1929, but in this analysis, we are only considering the 
span from 1944 to the present i.e. 2022 (disregarding the year 2023, because the data from this year is incomplete). 
We are not considering data from the start of the measurements due the fact that there are a lot of missing monthly records before this moment.

\subsection{Overview}%
\label{sub:overview}
The periodicity of the river flooding is significant. Visually, there is not much variation in the maximum discharge.

\begin{figure}[H]
    \centering
    \incfig{timeseries_full_overview}
    \caption{Full overview of the data considered in the analysis}
    \label{fig:full_mothly_overview}
 \end{figure}

Since the daily data is also available to us, we can observe the seasonality of discharges.

\begin{figure}[hbt]
    \centering
    \incfig{timeseries_groups_by_year}
    \caption{Seasonal overview of the discharge rate of the river determined from the daily measurements. The colour
        represents the year, when the measurements occurred.}
    \label{fig:seasonality_of_discharges}
 \end{figure}

\begin{table*}[h!]
    \centering
    \begin{tabularx}{\linewidth}{r>{\centering\arraybackslash}X>{\centering\arraybackslash}X>{\centering\arraybackslash}X>{\centering\arraybackslash}X>{\centering\arraybackslash}X>{\centering\arraybackslash}X>{\centering\arraybackslash}X}
\toprule  Month & $\bar{X}_n$ & $\std$ & $\min$ & $Q_1$ & $Q_2$ & $Q_3$ & $\max$\\\midrule
  January & \qty{1571.0}{\metre\cubed\per\second} & \qty{408.8}{\metre^6\per\second\squared} & \qty{295.0}{\metre\cubed\per\second} & \qty{1372.4}{\metre\cubed\per\second} & \qty{1701.3}{\metre\cubed\per\second} & \qty{1844.0}{\metre\cubed\per\second} & \qty{2365.0}{\metre\cubed\per\second}\\\addlinespace
  February & \qty{1283.2}{\metre\cubed\per\second} & \qty{601.8}{\metre^6\per\second\squared} & \qty{107.2}{\metre\cubed\per\second} & \qty{748.8}{\metre\cubed\per\second} & \qty{1421.6}{\metre\cubed\per\second} & \qty{1787.8}{\metre\cubed\per\second} & \qty{2365.0}{\metre\cubed\per\second}\\\addlinespace
  March & \qty{798.8}{\metre\cubed\per\second} & \qty{623.9}{\metre^6\per\second\squared} & \qty{32.4}{\metre\cubed\per\second} & \qty{251.4}{\metre\cubed\per\second} & \qty{603.4}{\metre\cubed\per\second} & \qty{1335.8}{\metre\cubed\per\second} & \qty{2251.2}{\metre\cubed\per\second}\\\addlinespace
  April & \qty{383.8}{\metre\cubed\per\second} & \qty{452.5}{\metre^6\per\second\squared} & \qty{9.0}{\metre\cubed\per\second} & \qty{88.4}{\metre\cubed\per\second} & \qty{180.0}{\metre\cubed\per\second} & \qty{460.9}{\metre\cubed\per\second} & \qty{1851.9}{\metre\cubed\per\second}\\\addlinespace
  May & \qty{151.9}{\metre\cubed\per\second} & \qty{211.9}{\metre^6\per\second\squared} & \qty{2.6}{\metre\cubed\per\second} & \qty{41.8}{\metre\cubed\per\second} & \qty{76.8}{\metre\cubed\per\second} & \qty{143.4}{\metre\cubed\per\second} & \qty{1342.6}{\metre\cubed\per\second}\\\addlinespace
  June & \qty{91.7}{\metre\cubed\per\second} & \qty{96.0}{\metre^6\per\second\squared} & \qty{0.0}{\metre\cubed\per\second} & \qty{25.2}{\metre\cubed\per\second} & \qty{62.2}{\metre\cubed\per\second} & \qty{118.8}{\metre\cubed\per\second} & \qty{759.0}{\metre\cubed\per\second}\\\addlinespace
  July & \qty{214.2}{\metre\cubed\per\second} & \qty{242.8}{\metre^6\per\second\squared} & \qty{0.2}{\metre\cubed\per\second} & \qty{50.8}{\metre\cubed\per\second} & \qty{115.1}{\metre\cubed\per\second} & \qty{292.0}{\metre\cubed\per\second} & \qty{1512.0}{\metre\cubed\per\second}\\\addlinespace
  August & \qty{732.1}{\metre\cubed\per\second} & \qty{468.5}{\metre^6\per\second\squared} & \qty{3.2}{\metre\cubed\per\second} & \qty{391.0}{\metre\cubed\per\second} & \qty{683.4}{\metre\cubed\per\second} & \qty{971.0}{\metre\cubed\per\second} & \qty{2500.9}{\metre\cubed\per\second}\\\addlinespace
  September & \qty{1192.3}{\metre\cubed\per\second} & \qty{339.5}{\metre^6\per\second\squared} & \qty{149.2}{\metre\cubed\per\second} & \qty{976.8}{\metre\cubed\per\second} & \qty{1141.4}{\metre\cubed\per\second} & \qty{1388.3}{\metre\cubed\per\second} & \qty{2658.0}{\metre\cubed\per\second}\\\addlinespace
  October & \qty{1259.9}{\metre\cubed\per\second} & \qty{160.4}{\metre^6\per\second\squared} & \qty{825.8}{\metre\cubed\per\second} & \qty{1137.7}{\metre\cubed\per\second} & \qty{1259.2}{\metre\cubed\per\second} & \qty{1381.7}{\metre\cubed\per\second} & \qty{1867.5}{\metre\cubed\per\second}\\\addlinespace
  November & \qty{1428.6}{\metre\cubed\per\second} & \qty{166.2}{\metre^6\per\second\squared} & \qty{997.9}{\metre\cubed\per\second} & \qty{1315.7}{\metre\cubed\per\second} & \qty{1434.5}{\metre\cubed\per\second} & \qty{1541.0}{\metre\cubed\per\second} & \qty{1997.3}{\metre\cubed\per\second}\\\addlinespace
  December & \qty{1585.1}{\metre\cubed\per\second} & \qty{217.4}{\metre^6\per\second\squared} & \qty{739.8}{\metre\cubed\per\second} & \qty{1463.5}{\metre\cubed\per\second} & \qty{1609.5}{\metre\cubed\per\second} & \qty{1725.6}{\metre\cubed\per\second} & \qty{2106.7}{\metre\cubed\per\second}\\\bottomrule
\end{tabularx}
    \caption{Elementary statistics of the data grouped by months}
    \label{tab:month_stats}
\end{table*}
  
 Before limiting ourselves to the maximum yearly discharges, let us look at the seasonality of the measurements. 
 For a comprehensive overview of elementary statistics of the seasonal data, see \Cref{tab:month_stats}.

 \begin{figure}[hbt]
     \centering
     \incfig{month_boxplot}
     \caption{Box plots of the monthly discharges}
     \label{fig:month_box_plots}
 \end{figure}

 From \Cref{fig:month_box_plots} we can see the same periodicity that was already apparent from
 \Cref{fig:full_mothly_overview,fig:seasonality_of_discharges}. We can also see that there are a lot of
 outliers in the data (in the sense of a Tukey style Box plot \cite{mcgillVariationsBoxPlots1978}, meaning they are
 greater than \(Q_3 + 1.5 \cdot \mathop{IQR}\)), which would suggest a heavy tailed distribution.

 \section{Extremal values}%
 \label{sec:extremal_values}
 From now on, we will consider the annual maxima and minima of discharge rates at the Niamey hydrology station. We will
 denote the maximum values \[
     R_{\max}(y) = \max_i \left\{ X^{(y)}_i \right\} \qcomma \tilde{R}_{\max}(y) = \max_{i \in [12]}\{ \tilde{X}^{(y)}_i \} 
 \] 
 and minimum values \[
     R_{\min}(y) = \min_i \left\{ X^{(y)}_i \right\} \qcomma \tilde{R}_{\min}(y) = \min_{i \in [12]}\{ \tilde{X}^{(y)}_i \}
 \] 
 where, \(\{X_{y}\}\) are the daily measurements (365-ish of them), \(\{\tilde{X}_{y}\}\) are the 12 monthly measurements 
 and \(y \in \left\{ x \in \N : 1944 \le x \le 2022 \right\}\) is the year of the measurement.

 \begin{figure}[h]
     \centering
     \incfig{R_max_R_min_timeseries}
     \caption{Comparison of yearly maxima and minima derived from the daily and monthly measurements}
     \label{fig:R_max_R_min_timeseries}
 \end{figure}

 From \Cref{fig:R_max_R_min_timeseries} we can see what is an obvious general rule, that the minima derived from the daily
 measurements are lower than those of monthly measurements and analogously, the maxima are higher than those
 derived from the monthly means. Because we are concerned mainly with flooding, the extrema derived from daily measurements are
 better suited for our analysis, due to the fact that a flood is not guaranteed to last for a full month and it does not
 care for our arbitrary month boundaries. We can argue similarly for daily measurements, but anyhow, they logically tend
 to smooth our data less and are therefore better suited for analysis of extremal values. Also for a comparison see
 \hyperref[fig:comp_max_min_day_month]{Figure \the\numexpr\thefigure+1\relax}.

 \begin{figure*}[h]
    \centering
   \begin{subfigure}{0.4\textwidth}
    \label{fig:comp_max_min_day_month}
    \centering
    \incfig{max_boxplot}
    \caption{Comparison of \(\tilde{R}_{\max}\) and \(R_{\max}\)}
   \end{subfigure}
   \begin{subfigure}{0.4\textwidth}
    \centering
    \incfig{min_boxplot}
    \caption{Comparison of \(\tilde{R}_{\min}\) and \(R_{\min}\)}
   \end{subfigure}
   \caption{Box plots for comparison of extremal values derived from daily and monthly measurements.}
 \end{figure*}

\begin{table*}[th]
        \centering
        \begin{tabularx}{\linewidth}{r>{\centering\arraybackslash}X>{\centering\arraybackslash}X>{\centering\arraybackslash}X>{\centering\arraybackslash}X>{\centering\arraybackslash}X>{\centering\arraybackslash}X>{\centering\arraybackslash}X}
\toprule  Variable & $\bar{X}_n$ & $\std$ & $\min$ & $Q_1$ & $Q_2$ & $Q_3$ & $\max$\\\midrule
  $\tilde{R}_{\max}$ & \qty{1744.0}{\metre\cubed\per\second} & \qty{268.4}{\metre^6\per\second\squared} & \qty{1105.8}{\metre\cubed\per\second} & \qty{1628.3}{\metre\cubed\per\second} & \qty{1792.6}{\metre\cubed\per\second} & \qty{1896.2}{\metre\cubed\per\second} & \qty{2315.5}{\metre\cubed\per\second}\\\addlinespace
  $R_{\max}$ & \qty{1856.9}{\metre\cubed\per\second} & \qty{291.5}{\metre^6\per\second\squared} & \qty{1153.9}{\metre\cubed\per\second} & \qty{1727.0}{\metre\cubed\per\second} & \qty{1875.7}{\metre\cubed\per\second} & \qty{2029.2}{\metre\cubed\per\second} & \qty{2658.0}{\metre\cubed\per\second}\\\addlinespace
  $R_{\min}$ & \qty{40.5}{\metre\cubed\per\second} & \qty{42.4}{\metre^6\per\second\squared} & \qty{0.0}{\metre\cubed\per\second} & \qty{6.0}{\metre\cubed\per\second} & \qty{31.4}{\metre\cubed\per\second} & \qty{54.4}{\metre\cubed\per\second} & \qty{211.9}{\metre\cubed\per\second}\\\addlinespace
  $\tilde{R}_{\min}$ & \qty{65.5}{\metre\cubed\per\second} & \qty{57.3}{\metre^6\per\second\squared} & \qty{3.3}{\metre\cubed\per\second} & \qty{23.0}{\metre\cubed\per\second} & \qty{60.0}{\metre\cubed\per\second} & \qty{93.1}{\metre\cubed\per\second} & \qty{300.9}{\metre\cubed\per\second}\\\bottomrule
\end{tabularx}
        \caption{Elementary statistics of extremal values}
        \label{tbl:max_stats}
\end{table*}


From now on, we will be working with the extremal variables derived from the daily observations i.e. \(R_{\max}\) and
\(R_{\min}\).

\subsection{Independence}%
\label{sub:independence}
For any kind of distribution-free estimation, we need to have the estimates of the mean \(\mu\). To approximate the
variance of the estimates, we first need to prove sequential independence of the time-series \(\left\{ R_{\max}(y) \right\}\) and 
\(\left\{ R_{\min}(y) \right\}\). As an alternative to testing independence, we will be testing sequential non-correlation.
That is, we are testing \[
    H_0: (\forall k \in \left[ h \right])(\rho_k(t) \coloneqq \Cov(X(t),\, X(t+k)) = 0),
\] 
where \(h\) is the maximum period of correlation that we allow and the hypothesis is taken for all \(t\) values that are
applicable. We are considering this hypothesis against the alternative of \[
    H_1: (\exists k \in \left[ h \right])(\rho_k(t) \neq 0).
\] 
The existence of such a \(k\) from hypothesis \(H_1\) would imply the time-series to have seasonality with the period of
\(k\) years.

We will be using the \emph{Ljung-Box test}. The statistic is \[
    Q = n(n+2) \sum_{k = 1}^{h} \frac{ \hat{\rho}_k^2 }{n - k}
\] and by \cite{ljungMeasureLackFit1978} under the assumption of \(H_0\), it holds asymptotically, that \(Q \sim \chi^2_{h}\).

\begin{table}[H]
  \centering
  \begin{tabularx}{\linewidth}{rccc}
\toprule  Test Name & $p$-value & $Q$ & $h$\\\midrule
  Ljung-Box autocorrelation test & 9.0e-11 & 49.8 & 3\\\addlinespace
  Box-Pierce autocorrelation test & 2.9e-10 & 47.4 & 3\\\bottomrule
\end{tabularx}
  \caption{Ljung-Box test of serial correlation (independence) and Box-Pierce test for completeness, even though it
  is a only a simplified version of the Ljung-Box test.}
  \label{tab:indep_test}
\end{table}

By the tests in \Cref{tab:indep_test}, we have to conclude that the time-series is not serially independent up to the lag
of \(3 \cdot\Delta t\) (in this table \(\Delta t = 1\) year which evaluates to a lag of 3 years). However, we can increase the delays between the samples and test again for this new, smaller sample.

\begin{table}[H]
  \centering
  \begin{tabularx}{\linewidth}{>{\centering\arraybackslash}X>{\centering\arraybackslash}X>{\centering\arraybackslash}X>{\centering\arraybackslash}X>{\centering\arraybackslash}X}
\toprule  $p$-value & $Q$ & $n$ & $h$ & $\Delta t$\\\midrule
  0.077 & 6.84 & 40 & 3 & 2 years\\\addlinespace
  0.015 & 10.5 & 27 & 3 & 3 years\\\addlinespace
  0.38 & 3.1 & 20 & 3 & 4 years\\\bottomrule
\end{tabularx}
  \caption{Ljung-Box autocorrelation test for increased delay between the samples}
  \label{tab:indep_test_wide}
\end{table}

As we can see from tests in \Cref{tab:indep_test_wide}, we would eventually get to a delay for which the 
correlation would be minute enough to state that there is no reason to suggest serial dependence of the time-series.


Unfortunately, our original sample has only \(79\) years worth of observations, so the assurance of independence would not 
be a sufficient reason to make this compromise on the asymptotics (arising from the loss of considered measurements). We will keep the full dataset and continue with the
grain of salt that our estimates may not have a theoretically proven asymptotic precision. 

\subsection{stationarity}%
\label{sub:stationarity}

One property that we will need for the estimates to be relevant however is strong stationarity. In other words, we need to
have confidence that the data comes from a single distribution, otherwise there is no reason for even estimating it. In
the terms of our data, we want to test that there is no apparent trend. That is, we would like to rule out the
possibility of water in the Niger river decreasing or increasing with time due to some external factor (e.g. climate
change). A great fit for this kind of statistical test would be \emph{Page's trend test} from \cite{pageContinuousInspectionSchemes1954}.

However I will only intuitively test weak stationarity by testing that even samples have the same mean as the odd 
samples using \emph{ANOVA} and similarly using the \emph{Levene test} for equal variances.

\begin{table}[H]
  \centering
  \begin{tabularx}{\linewidth}{r>{\centering\arraybackslash}X>{\centering\arraybackslash}X}
\toprule  Test Name & $p$-value & $F$\\\midrule
  One-way ANOVA & 0.27 & 1.25\\\bottomrule
\end{tabularx}
  \caption{One way Analysis of Variance test for equality of means applied to even and odd samples of the time-series}
  \label{tab:mean_test}
\end{table}

From the test in \Cref{tab:mean_test} we will further assume equal means.

\begin{table}[H]
  \begin{tabularx}{\linewidth}{r>{\centering\arraybackslash}X>{\centering\arraybackslash}X}
\toprule  Test Name & $p$-value & $W$\\\midrule
  Levene's test & 0.97 & 0.00155\\\bottomrule
\end{tabularx}
  \caption{Levene's test of equality of variances centred by the sample means and applied to even and odd samples of the
  time-series.}
  \label{tab:var_test}
\end{table}

From the test in \Cref{tab:var_test} we will assume equal variances and without further testing assume strong
stationarity. We can now move on to the distribution-free estimation with an indication that the data comes from one
distribution and therefore the estimates will make sense.

\section{Objective Values}%
\label{sec:objective_values}

For making the estimates in the objectives, we first need to determine the values of the objectives.
For \(C\) we have
\[
    \begin{split}
        C &= \frac{Q_3 +  (R_{\max})_{[3n/4] + 1}}{2} = \\
          &= \frac{\qty{2029.225}{\metre\cubed\per\second} + \qty{2033.78}{\metre\cubed\per\second}}{2} = \\
          &= \qty{2031.5025}{\metre\cubed\per\second}
    \end{split}
\]
for \(D\) we have
\[
    D = (R_{\max})_n + \qty{100}{\metre\cubed\per\second} = \qty{2757.965}{\metre\cubed\per\second}.
\]
From \Cref{obj:III}, we have \[
    \begin{split}
        E &= \frac{(R_{\max})_{4} +  (R_{\max})_{5}}{2} = \\
          &= \frac{\qty{1266.13}{\metre\cubed\per\second} + \qty{1282.64}{\metre\cubed\per\second}}{2} = \\
          &= \qty{1274.385}{\metre\cubed\per\second}.
    \end{split}
\] 
and for \(F\) (denoting the value in \Cref{obj:V}) we have
\[
    F = (R_{\max})_n = \qty{2657.965}{\metre\cubed\per\second}.
\] 
From \Cref{obj:VI}, we also need to know the value for \(Z\) which is \[
    Z = \frac{(R_{\max})_n - Q_3}{2} = \qty{314.37}{\metre\cubed\per\second}.
\] 

\begin{table}[h!]
    \centering
    \begin{tabularx}{\linewidth}{r>{\centering\arraybackslash}X>{\centering\arraybackslash}X}
\toprule  Variable & value & Formula\\\midrule
  $C$ & $\qty{2031.502}{\metre\cubed\per\second}$ & $\frac{Q_3 +  (R_{\max})_{[3n/4] + 1}}{2}$\\\addlinespace
  $D$ & $\qty{2757.965}{\metre\cubed\per\second}$ & $(R_{\max})_n + \qty{100}{\metre\cubed\per\second}$\\\addlinespace
  $E$ & $\qty{1274.385}{\metre\cubed\per\second}$ & $\frac{(R_{\max})_{4} +  (R_{\max})_{5}}{2}$\\\addlinespace
  $F$ & $\qty{2657.965}{\metre\cubed\per\second}$ & $(R_{\max})_n$\\\addlinespace
  $Z$ & $\qty{314.37}{\metre\cubed\per\second}$ & $\frac{(R_{\max})_n - Q_3}{2}$\\\bottomrule
\end{tabularx}
    \caption{Values that occur in the objectives}
\end{table}


\section{Distribution-free Estimates}%
\label{sec:distribution-free_estimates}

We can make distribution-free estimates for objectives \Cref{obj:I,obj:II,obj:V}.

We will calculate an upper bound of the probability using Markov's inequality \[
    P(X > \varepsilon) \le \mathbb{E}X/\varepsilon \eqqcolon MUB,  
\] and Chebyšev's-Cantelli's inequality \[
    P(X - \mathbb{E}X > \varepsilon) \le \frac{\mathbb{D}X}{\mathbb{D}X + \varepsilon^2},
\]  therefore \[
     P(X > \varepsilon + \mathbb{E}X) \le \frac{\mathbb{D}X}{\mathbb{D}X + \varepsilon^2} \eqqcolon CCUB.
\] 
\begin{table}[H]
  \begin{tabularx}{\linewidth}{r>{\centering\arraybackslash}X>{\centering\arraybackslash}X}
\toprule  Threshold & $MUB$ & $CCUB$\\\midrule
  $C$ & 0.914 & 0.736\\\addlinespace
  $D$ & 0.673 & 0.0947\\\addlinespace
  $F$ & 0.699 & 0.117\\\bottomrule
\end{tabularx}
  \caption{Distribution-free estimates for \Cref{obj:I,obj:II,obj:V}}
  \label{tab:dist_free}
\end{table}

We also have a distribution-free estimate for the lower bound in the \(L^{p}\) version of 
\emph{Paley-Zygmund inequality} \cite{petrovLowerBoundsTail2007} by
\[
    P(X > \theta\cdot \mathbb{E}X) \ge \frac{(1-\theta)^{p/(p-1)} \cdot \mathbb{E}[X]^{p/(p-1)}}{\mathbb{E}[X^p]^{1/(p-1)}} \eqqcolon PZLB_{p},
\] which holds for \(\theta \in \left< 0, 1 \right>\). This makes it a bad fit for the lower bounds of all of \(C\), \(D\)
and \(F\), since
\[
    \theta_{D} \ge \theta_{F} \ge \theta_{C} = C/\mathbb{E}[R_{\max}] \approx \frac{\qty{2031.5025}{\metre\cubed\per\second}}{\qty{1856.9178}{\metre\cubed\per\second}}
    \approx 1.095,
\] and this lower bound estimate is not good for \(\theta\) near \(0\) and \(1\) anyway.

But we can use this for
estimating the upper bound of the drought level in \Cref{obj:III}, since
\begin{equation}
    \label{eq:PZLB_p}
    P(X \le E) = 1 - P(X > E) \le 1 - PZLB_{p}(E).
\end{equation}
Now from \Cref{eq:PZLB_p} we have an upper bound on the probability that is a function of \(p\), so we can also
maximize the \(PZLB_{p}\) in the \(p\) variable leading to the optimal upper bound constructed in this manner of \[
    P(R_{\max} \le E) \le 1 - PZLB_{12.7}(E) \approx 1 - 0.25104 = 0.74896,
\] and that is our distribution-free estimate for the upper bound of the drought probability, which unfortunately is
a very imprecise bound.

\begin{table}[h!]
    \centering
    \begin{tabularx}{\linewidth}{r>{\centering\arraybackslash}X}
\toprule  Objective & Estimate\\\midrule
  \ref{obj:I} & $P[R_{\max} > C] \le 0.74$\\\addlinespace
  \ref{obj:II} & $P[R_{\max} > D] \le 0.095$\\\addlinespace
  \ref{obj:V} & $P[R_{\max} > F] \le 0.12$\\\addlinespace
  \ref{obj:III} & $P[R_{\max} \le E] \le 0.75$\\\bottomrule
\end{tabularx}
    \caption{Summary of the distribution free estimates for the objectives}
    \label{tab:dist_free_summary}
\end{table}


\section{Non-parametric Distribution Estimates}%
\label{sec:non-parametric_estimates}

We want to estimate the probability (not the bounds) for all of the objectives. For this we can also employ
non-parametric estimates. One such way of estimation is using a non-parametric density estimate and determining the bound
from this.

\subsection{Histogram Density Estimate}%
\label{sub:histogram_density_estimate}

The histogram estimate of the probability density of the random variable \(X\), for the (left closed) bins \[
    \begin{split}
        B_1 &= \left[\min(X)\eqqcolon m, m + \frac{M - m}{k}\right),\\
        B_2 &= \left[m + \frac{ M - m }{k}, m + 2\cdot\frac{M - m}{k}\right),\\
            &\vdots\\
        B_k &= \left[m + (k-1)\cdot\frac{M - m}{k}, M \coloneqq \max(X)\right]\\
    \end{split}
\]  is \begin{equation}
    \label{eq:histogram_density_estimate}
    \hat{p}_n^{H}(x) = \frac{k}{n}\sum_{i = 1}^{n} I(X_i \in B_\ell), \qq{such that} x \in B_\ell.
\end{equation}

We will use various methods for determining the optimal bin-size \(k\).
This means we are interested in selecting an optimal \(h\), where the number of bins is \[
    k = \left\lceil \frac{\max(R_{\max}) - \min(R_{\max})}{h} \right\rceil.
\] 

\subsubsection{Square-root choice}%
\label{ssub:square-root_choice}

Suggested bin-count \[
    k = \left\lceil \sqrt{n} \right\rceil.
\]
For our data this is \[
    k = \left\lceil \sqrt{79} \right\rceil = 9.
\] 

\subsubsection{Struge's Formula}%
\label{ssub:struge's_formula}

This method is derived from the binomial distribution and it suggests \[
    k = \left\lceil \log_2(n) \right\rceil + 1.
\]
For our data this is \[
    k = \left\lceil \log_2(79) \right\rceil + 1 = 7.
\] 

\subsubsection{Rice's Formula}%
\label{ssub:rice_formula}

Suggested bin-count \[
    k = \left\lceil 2 \sqrt[3]{n} \right\rceil.
\] 
For our data this is  \[
    k = \left\lceil 2 \sqrt[3]{79} \right\rceil = 9.
\] 


\subsubsection{Doane's Formula}%
\label{ssub:doane's_formula}

Suggested bin-count \[
    k = 1 + \left\lceil \log_2(n) + 
    \log_2 \left( 1 + \frac{\left| \widehat{\tilde{\mu}_3} \right|}{\sigma_{\widehat{\tilde{\mu}_3}}} \right)
\right\rceil,
\] where \(\widehat{\tilde{\mu}_3}\) is an estimator of skewness and \(\sigma_{\widehat{\tilde{\mu}_3}}\) is the
\ac{SD} of the estimate. For the estimator \[
    g_1 = \frac{m_3}{m_2^{3/2}}\qcomma m_2 = \widehat{\mathbb{E}[R_{\max}^2]} \qqtext{and} m_3
    = \widehat{\mathbb{E}[R_{\max}^3]},
\] the standard deviation is \[
    \sigma_{g_1} = \sqrt{\frac{6(n-2)}{(n+1)(n+3)}}.
\] 
For our data this suggests the bin-count of \[
    k = 9.
\]

\subsubsection{Freedman-Diaconi's Choice}%
\label{ssub:freedman-diaconi's_choice}

This rule suggests \[
    h = 2 \frac{\mathop{IQR}(x)}{\sqrt[3]{n}},
\] which for our data equates to \[
h \approx \qty{140.85}{\metre\cubed\per\second}.
\] and \[
    k = 11.
\] 

\begin{table*}[t]
  \begin{tabularx}{\linewidth}{ll>{\centering\arraybackslash}X>{\centering\arraybackslash}X}
\toprule  Method & Formula & $k$ & $h$\\\midrule
  Square-root & $k = \left\lceil \sqrt{n} \right\rceil$ & 9 & \qty{167.12}{\metre\cubed\per\second}\\\addlinespace
  Sturge & $k = \left\lceil \log_2(n) \right\rceil + 1$ & 8 & \qty{188.01}{\metre\cubed\per\second}\\\addlinespace
  Rice & $k = \left\lceil 2 \sqrt[3]{n} \right\rceil$ & 9 & \qty{167.12}{\metre\cubed\per\second}\\\addlinespace
  Doan & $k = 1 + \left\lceil \log_2(n) + \log_2 \left( 1 + \frac{\left| g_1 \right|}{\sigma_{g_1}}\right) \right\rceil$ & 9 & \qty{167.12}{\metre\cubed\per\second}\\\addlinespace
  Freedman-Diaconi & $h = 2 \frac{\mathop{IQR}(x)}{\sqrt[3]{n}}$ & 11 & \qty{140.86}{\metre\cubed\per\second}\\\bottomrule
\end{tabularx}
  \caption{Summary of methods for histogram parameter selection}
  \label{tab:hist_choices}
\end{table*}

\begin{figure}[H]
    \centering
  \begin{subfigure}{0.4\textwidth}
    \centering
    \incfig{histogram_with_8_bins}
  \end{subfigure}

  \begin{subfigure}{0.4\textwidth}
    \centering
    \incfig{histogram_with_9_bins}
  \end{subfigure}
  \begin{subfigure}{0.4\textwidth}
    \centering
    \incfig{histogram_with_11_bins}
  \end{subfigure}
  \caption{Histograms of \(R_{\max}\) for various bin-counts}
  \label{fig:histograms}
\end{figure}


\subsubsection{Histogram Estimates}%
\label{ssub:histogram_estimates}
We will select the histogram with \(11\) bins suggested by the choice \ref{ssub:freedman-diaconi's_choice}, because
the estimate of the \ac{IQR} is robust even if our data sample is not independent. We will be able to make
estimations for the threshold of \Cref{obj:I} and for the probability \(0.1\) of \Cref{obj:IV,obj:VI}.

Estimating the probability of the discharge passing the level \(T\) using a histogram estimate means calculating the
definite integral from \(T\) to \(+\infty\) of the histogram estimate of the probability density function at
\Cref{eq:histogram_density_estimate}. That is \[
    \hat{P}[X > T] = \int_{T}^{+\infty} \hat{p}_n^{\text{H}}(x) \dd{x}.
\] 

The reason why we need to be careful about this estimate and can not use it blindly for estimating any probability is
that if the threshold \(T\) is close to the extrema of the data, the estimated \(\hat{p}^H_n(x)\) in the integrand 
may be under or over estimated (has a large standard deviation) due to the lack of measurements in the integration
region. At the right edge of the histogram this is even more so
for distributions that are heavy-tailed. We also do not have an estimate of the probability density at values that are
greater than the maximum from the data or lower than the minimum of it and we cannot be sure, where the support of the
density really is. Moreover, using this definition of a histogram (one that is normalized to sum up to 1) does not 
allow us to estimate the probability of passing any threshold that is larger than the maximum of the data and that
limits us to the decade flood threshold (i.e. eliminating the ability to estimate for the probabilities of \(0.04\),
\(0.02\), \(0.01\) and \(0.004\) of objectives \Cref{obj:IV,obj:VI}).


That is the reason why we can only estimate these selected objectives from the histogram and not the rest, since they
are not represented in the data. For example \Cref{obj:III} cannot be estimated from the histogram we have here,
because, we do not know what the probability density estimate is under the minimum (it is equal to \(0\) using 
the histogram estimate) of the observed data and thus our estimate would not take this under consideration. We
will use the \ac{KDE} for estimating the other objectives. 


For \Cref{obj:I} we have the estimate of \[
    \hat{P}\left[R_{\max} \ge C\right] = \int_{C}^{+\infty} p_n^{H}(x) \dd{x} \approx 0.257.
\] 

Similarly for the threshold in \Cref{obj:IV} of the probability \(0.1\) we have \[
    \hat{P}^{-1}(0.1) = \qty{2181.665}{\metre\cubed\per\second} \eqqcolon T, 
\] where \[
    \hat{P}(T) = \hat{P}[ R_{\max} \ge T ].
\] 

For \Cref{obj:VI} we have to be careful about the question. It is asking about the conditional probability. The
estimate is as follows \[
    \begin{split}
        \hat{P}[ Z \le R_{\max} - C \mid R_{\max} \ge C ] &= \frac{\hat{P}[ Z + C \le R_{\max} ]}{\hat{P}[ R_{\max} \ge C ]}\\
                                                        &\approx \frac{0.045}{0.257} \approx 0.175.
    \end{split}
\] 

\begin{table}[h!]
    \centering
    \begin{tabularx}{\linewidth}{r>{\centering\arraybackslash}X}
\toprule  Objective & Estimate\\\midrule
  \ref{obj:I} & $\hat{P}[R_{\max} > C] = 0.257$\\\addlinespace
  \ref{obj:IV} & $\hat{h}_{0.1} = \qty{2181.67}{\metre\cubed\per\second}$\\\addlinespace
  \ref{obj:V} & $\hat{P}[R_{\max} \ge Z + C \mid R_{\max} \ge C] = 0.176$\\\bottomrule
\end{tabularx}
    \caption{Summary of the non-parametric histogram estimates for the objectives}
    \label{tab:histogram_estimates_summary}
\end{table}

A summary of the estimates can be seen in \Cref{tab:histogram_estimates_summary}.


\subsection{Kernel Density Estimate}%
\label{sub:kernel_density_estimate}

\ac{KDE} is a non-parametric method for estimating the probability density function of a random variable. It works by assigning a weight to each data point, and then summing the weighted kernels over the range of the variable. The kernel function is typically a bell-shaped curve, such as the Gaussian.

The following equation shows the \ac{KDE} of a random variable $X$ with data points $x_1, x_2, ..., x_n$:

\[
    \hat{p}_n^{\text{K}}(x) = \frac{1}{n} \sum_{i=1}^n K\left( \frac{x - x_i}{\sigma} \right),
\]

where \(K\) is the kernel function, \(\sigma\) is the bandwidth parameter, and \(\hat{p}_n^{\text{K}}(x)\) is the estimated probability density function.

We will be choosing the bandwidth parameter and the kernel function visually for our data sample. The visual cues that
will determine our choice are the following:
\begin{description}
    \item[\label{itm:desc_cues_smoothness}Smoothness] \(\hat{p}_n^{\text{K}}\) should be smooth, not jagged and should not have sharp edges
    \item[\label{itm:desc_cues_consistency}Consistency] \(\hat{p}_n^{\text{K}}\) should be visually similar to the histogram estimate \(\hat{p}_n^{\text{H}}\)
\end{description}

The accuracy of the \ac{KDE} depends more no the choice of of \(\sigma\) than on the kernel \(K\) \cite{markovichClassicalMethodsProbability2007}
and therefore, we will be concentrating on the choice of the bandwidth parameter and use the, perhaps most widely used,
Epanechnikov kernel
\[
    K(x)=\frac{3}{4}(1-x^{2}).
\] 
Furthermore, we have extremal data, meaning that the distribution is presumably heavy-tailed. This creates a problem
of under-representation of the probability in the tail. There are three conventional methods for fighting this problem.

\begin{enumerate}[label=\alph*)]
    \item combined parametric non-parametric methods
    \item variable bandwidth kernel estimator
    \item retransformed non-parametric estimator
\end{enumerate}

I will concentrate on the third option because of its simplicity and the possibility not to make any assumptions about
the distribution family. We will transform the data using the transformation

\begin{align}
    T(x) &= (2/\pi)\cdot \arctan(x) \\
    T^{-1}(y) &= \tan((\pi/2)\cdot y).
\end{align}

and the procedure is to transform the data creating a new \ac{RV} \(Y\), perform the non-parametric density
estimate and then use the inverse transform on the estimate produced. The point in this is that the support of the
transformed \ac{RV} is compact and therefore more suitable for \ac{KDE}
\cite{markovitchNonparametricEstimationLongtailed2000, markovichClassicalMethodsProbability2007}.

A good starting point for the bandwidth parameter \(\sigma\) is the \emph{Silverman's rule of thumb}
\cite{silvermanDensityEstimationStatistics1986}.
\[
    \sigma_{\text{Silverman's}} = 0.9 \cdot \min \left\{  \hat{\sigma}, \frac{\mathop{IQR}}{1.34} \right\} \cdot n^{-\frac{1}{5}},
\] 

where \(\mathop{IQR}\) is the interquartile range and \(\hat{\sigma}\) is the standard deviation estimate. This is
a rule choice that is optimal for a normally distributed sample. We will use is as a starting point for our choice.

The Silverman's rule of thumb gives 
\[
    \sigma_{\text{Silverman's}} \approx 1.538\times 10^{-5}
\]
for our transformed data \(T(X)\). We will make our choice visually considering the neighbourhood of the choice 
\(\sigma_{\text{Silverman's}}\) with the comparison with the histogram and taking into consideration
\hyperref[itm:desc_cues_consistency]{consistency} in the \Cref{fig:kde_choice_lines}.

\begin{figure}[H]
    \centering
    \incfig{KDE_comparison_lines}
    \caption{Comparison of different values of the bandwidth parameter from the neighbourhood of the Silverman's rule of
    thumb value with an overlay of the selected histogram.}
    \label{fig:kde_choice_lines}
\end{figure}

We can also visualize the influence of the bandwidth parameter \(\sigma\) on \(\hat{p}^{\text{K}}_{n}\) using
a surface plot with the neighbourhood around the initial choice \(\sigma_{\text{Silverman's}}\) in figure
\Cref{fig:kde_choice_surface}.

\begin{figure}[H]
    \centering
    \includegraphics[width=0.8\columnwidth]{KDE_comparison_surface.png}
    \caption{Comparison of different values of the bandwidth parameter from the neighbourhood of the Silverman's rule of
    thumb value with an overlay of the selected histogram. The highlighted lines are the estimates with the selected bandwidth parameter.}
    \label{fig:kde_choice_surface}
\end{figure}

With consideration of the criteria of \hyperref[itm:desc_cues_consistency]{consistency} and \hyperref[itm:desc_cues_smoothness]{smoothness} in the
\Cref{fig:kde_choice_surface} a good compromise is the value of \(\sigma = 1.7\cdot 10^{-5}\). The final density estimate 
\(\hat{p}^{\text{K}}_{n}\) is shown in \Cref{fig:kde_choice_plane}.

\begin{figure}[H]
    \centering
    \incfig{KDE_final_plane}
    \caption{\ac{KDE} with the selected bandwidth parameter \(\sigma = 1.7\cdot 10^{-5}\).}
    \label{fig:kde_choice_plane}
\end{figure}

\subsubsection{Kernel Density Objective Estimates}%
\label{ssub:kernel_density_objective_estimates}

Now that we have the \ac{KDE}, we can perform the estimates for \Cref{obj:I,obj:II,obj:III,obj:IV,obj:V,obj:VI}.

For \Cref{obj:I,obj:II,obj:V} we have the estimate
\[
    \hat{P}[ X > T ] = \int_{T}^{+\infty} \hat{p}_{n}^{\text{K}}(x) \dd{x}.
\]

\begin{table}[h!]
    \centering
    \begin{tabularx}{\linewidth}{r>{\centering\arraybackslash}X}
\toprule  Objective & Estimate\\\midrule
  \ref{obj:I} & $\hat{P}[R_{\max} > C] = 0.26$\\\addlinespace
  \ref{obj:II} & $\hat{P}[R_{\max} > D] = 0.0053$\\\addlinespace
  \ref{obj:V} & $\hat{P}[R_{\max} > F] = 0.011$\\\bottomrule
\end{tabularx}
    % \caption{}
\end{table}

For \Cref{obj:III} we have
\[
     \hat{P}[ X < T ] = \int^{T}_{0} \hat{p}_{n}^{\text{K}}(x) \dd{x}
 \] and similarly for \Cref{obj:VI}.

\begin{table}[h!]
    \centering
    \begin{tabularx}{\linewidth}{r>{\centering\arraybackslash}X}
\toprule  Objective & Estimate\\\midrule
  \ref{obj:III} & $\hat{P}[R_{\max} < E] = 0.052$\\\addlinespace
  \ref{obj:VI} & $\hat{P}[R_{\max} \ge Z + C \mid R_{\max} \ge C] = 0.2$\\\bottomrule
\end{tabularx}
    % \caption{}
\end{table}

For finding the thresholds \(h_{p}\) in \Cref{obj:IV}, it holds that
\[
    \hat{P}[ X > h_{p} ] = p,
\]
which translates into finding the root of the function 
\[
    f(T) = \hat{P}[ X > T ] - p = \int_{T}^{+\infty} \hat{p}_{n}^{\text{K}}(x) \dd{x} - p.
\] 


\begin{table}[h!]
    \centering
    \begin{tabularx}{\linewidth}{r>{\centering\arraybackslash}X>{\centering\arraybackslash}X}
\toprule  \(p\) & Name & Estimate\\\midrule
  0.1 & 10-year water & $\hat{h}_{0.1} = \qty{2218.9}{\metre\cubed\per\second}$\\\addlinespace
  0.04 & 25-year water & $\hat{h}_{0.04} = \qty{2403.2}{\metre\cubed\per\second}$\\\addlinespace
  0.02 & 50-year water & $\hat{h}_{0.02} = \qty{2549.0}{\metre\cubed\per\second}$\\\addlinespace
  0.01 & 100-year water & $\hat{h}_{0.01} = \qty{2665.8}{\metre\cubed\per\second}$\\\addlinespace
  0.004 & 250-year water & $\hat{h}_{0.004} = \qty{2797.2}{\metre\cubed\per\second}$\\\bottomrule
\end{tabularx}
    % \caption{}
\end{table}

\section{Parametric Estimates}%
\label{sec:parametric_estimates}

For a parametric estimate, we first need to find a parametric distribution that we expect the data comes from.

\subsection{Distribution Selection}%
\label{sub:distribution_selection}

We want to create \ac{MLE} fits of different distributions to the data while ensuring consistency with the previous
\ac{KDE} and histogram estimate of the density. From the previous estimates (e.g. \Cref{fig:kde_choice_plane}), it seems
apparent that the distribution is bimodal (explaining why the hypotheses that we simplified our requirement of
stationarity to in \Cref{sub:stationarity} didn't yield the desired result). As most of our hypothetical distributions
are unimodal, it is nonsensical to test these without filtering the data first (i.e. selecting data only from one of the
modes), but even with this precaution, the estimates would be appropriate only under the assumption that the modes
are isolated (i.e. that the distribution from the other mode does not affect the one that we would hypothetically be
fitting in isolation).

This models our physical expectation of the river discharge process well since, it would hypothetically arise from the
physical model of the river if there are two modes of the rivers discharge operation determined by external forces (e.g.
temperature, climate, etc.). This means that determined by these factors, the discharge in a particular year could be
modelled by a single one of the distributions from the mixture.

A mixture model complicates our hypothesised distribution however, because we are forced to consider mixtures of
different unimodal ones (extending the number of possibilities from \(n\) to \(n^2\)). After preliminary fits of
combinations of \emph{Gumbel}, \emph{Gamma} and \emph{Weibull} distribution mixtures, I selected two parametric
distributions.\footnote{The likelihood equations system is not guaranteed to be convex for mixture
distributions in particular, so it is necessary to find initial parameters \(p_0\) for the distribution that are
close enough to the optimal solution to ensure that the \ac{MLE} is the global solution.}
Namely a mixture containing two Gumbel distributions and one that explains the first mode with a Gamma distribution and
the second with a Gumbel. The two Q-Q plots of the parametric fits can be seen in \Cref{fig:QQ_selection_mixture} and we
can compare them with the previous \ac{KDE} density function in \Cref{fig:PD_selection_mixture}.


\begin{figure}[H]
    \centering
    \incfig{parametric_QQ_selection}
    \caption{Q-Q plots of the two selected parametric distribution models.}
    \label{fig:QQ_selection_mixture}
\end{figure}

\begin{figure}[H]
    \centering
    \incfig{parametric_PD_selection}
    \caption{Probability density function plots of the two selected parametric distribution models.}
    \label{fig:PD_selection_mixture}
\end{figure}

% Due to this, we will consider
% only the most expected distributions in the mixtures, which we will assume are the ones that occur in river discharge
% data commonly where the discharges are from a single mode distribution and including ones that very flexible. I will
% include in this the \emph{Gumbel}, \emph{Gamma} and \emph{Weibull} distributions . This gives \(9\) possibilities for
% the distribution of our data.

\subsection{Estimates}%
\label{sub:estimates}

I will be selecting from the distribution as a mixture of two Gumbel distributions based on the Q-Q plot in 
\Cref{fig:QQ_selection_mixture} and based on the assumption that the distributions model of the modes are affected by
the properties of the river basin and shape as opposed to it being affected by the external forces. We can also see
how close the model could be to the \ac{KDE} density comparing the \ac{MLE} density with the parametric model that
minimizes the \(\ell_2\) norm of the difference of \(\hat{p}_n^{\text{K}}\) and the model \(\hat{p}_{\ell_2}\) in
\Cref{fig:gumbel_mixture_fits}.

% \begin{figure*}[h!]
%     \centering
%     \incfigwide{parametric_gumbel_mixture}
%     \caption{Fits of the mixture of two Gumbel distributions maximizing the likelihood resp. minimizing the  \(\ell_2\) norm.}
%     \label{fig:gumbel_mixture_fits}
% \end{figure*}

\begin{figure}[h]
    \centering
    \incfig{parametric_gumbel_mixture_vert}
    \caption{Fits of the mixture of two Gumbel distributions maximizing the likelihood resp. minimizing the  \(\ell_2\) norm.}
    \label{fig:gumbel_mixture_fits}
\end{figure}

The objective estimates are made in the same way as are made from the previous density estimates.

\[
    \mathop{CCDF}(T) \coloneqq \hat{P}[ X > T ] = \int_{T}^{+\infty} \hat{p}_{\text{MLE}}(x) \dd{x}.
\]

\begin{table}[h!]
    \centering
    \begin{tabularx}{\linewidth}{r>{\centering\arraybackslash}X}
\toprule  Objective & Estimate\\\midrule
  \ref{obj:I} & $\hat{P}[R_{\max} > C] = 0.23$\\\addlinespace
  \ref{obj:II} & $\hat{P}[R_{\max} > D] = 0.0026$\\\addlinespace
  \ref{obj:V} & $\hat{P}[R_{\max} > F] = 0.005$\\\bottomrule
\end{tabularx}
    % \caption{}
\end{table}

In \Cref{obj:III} we use
\[
    \mathop{CDF}(T) \coloneqq \hat{P}[ X < T ] = \int^{T}_{0} \hat{p}_{\text{MLE}}(x) \dd{x}
\]

and for the estimate in \Cref{obj:VI}, we have
\[
\hat{P}[ Z \le R_{\max} - C \mid R_{\max} \ge C ] = \frac{\hat{P}[ Z + C \le R_{\max} ]}{\hat{P}[ R_{\max} \ge C ]}.
\] 

\begin{table}[h!]
    \centering
    \input{tables/parametric_estimates_B.tex}
    % \caption{}
\end{table}


For the thresholds in \Cref{obj:IV} we can use the quantile function of the estimated model.
\begin{table}[h!]
    \centering
    \begin{tabularx}{\linewidth}{r>{\centering\arraybackslash}X>{\centering\arraybackslash}X}
\toprule  \(p\) & Name & Estimate\\\midrule
  0.1 & 10-year water & $\hat{h}_{0.1} = \qty{2175.9}{\metre\cubed\per\second}$\\\addlinespace
  0.04 & 25-year water & $\hat{h}_{0.04} = \qty{2326.2}{\metre\cubed\per\second}$\\\addlinespace
  0.02 & 50-year water & $\hat{h}_{0.02} = \qty{2437.4}{\metre\cubed\per\second}$\\\addlinespace
  0.01 & 100-year water & $\hat{h}_{0.01} = \qty{2547.7}{\metre\cubed\per\second}$\\\addlinespace
  0.004 & 250-year water & $\hat{h}_{0.004} = \qty{2692.8}{\metre\cubed\per\second}$\\\bottomrule
\end{tabularx}
    % \caption{}
\end{table}

\section{The Number of Record Values}%
\label{sec:the_number_of_record_values}%

\Cref{obj:V} also requires an estimate for the number of record values that will occur in the next 50 years. For
a sequence of \(n\) \ac{iid} random variables \((X_i)_{i=1}^n\), we define a maximal record indicator as \[
    Y_i = \begin{cases}
        1 & \text{if } i = 1 \text{ or } i > 1 \text{ and } X_i > \max_{j < i} \{ X_j\}\\
        0 & \text{otherwise.}
    \end{cases}
\] 

With this notation, what we consider the number of records in \(n\) time periods is \[
    M_n = \sum_{i = 1}^{n} Y_i.
\] 

What allows us to determine the probability \(P[Y_i = 1]\) is the realization that \[
    X_i > \max_{j < i} \{ X_j \} \Leftrightarrow F(X_i) > \max_{j < i} \{ F(X_j)\}),
\] 
where \(F(X)\) is the cumulative distribution function of \(X\). We know that \(F(X)\) has a uniform distribution
\(U_{[0,1]}\) and thus for the maximal record value indicator (an analogous derivation can be done for the minimum record
value indicator with the same result) it holds that \[
    P[Y_i = 1] = 1/i.
\] 

Thus the expected number of records in \(n\) time periods~is \[
    \mathbb{E}[M_n] = \mathbb{E}\left[ \sum_{i = 1}^{n} Y_i \right] = \sum_{i = 1}^{n} 1/i.
\] 

This however further argues in the direction of our sequence not being \ac{iid}, because the expected number of records
in our sequence of measurements would be \(\sum_{i = 1}^{79} 1/i \approx 5\) but there are \(13\) record values in our
sample. There is a test for assessing whether a sample is \ac{iid} based on this result, that is derived in
\cite{benestadRecordvaluesNonstationarityTests2004} that is based on determining the confidence interval using Monte
Carlo simulations. Using a similar simulation the \ac{SD} of \(M_n\) with \(n \equiv 79\) is \(1.8217\). The reliability
of the sample \(M_n\) being different from the theoretical and therefore the sample not being \ac{iid} is
\(4.42\sigma\).

However if we considered an \ac{iid} sample, then the number of expected record measurements in the next \(50\) years
(i.e. from 2023 to 2073) is \[
    M_{[80,129]} \coloneqq \mathbb{E}\left[ \sum_{i = 80}^{129} Y_i \right] = \sum_{i = 80}^{129} 1/i \approx 0.49.
\] 
A Monte Carlo simulation of the \ac{SD} of \(M_{[80,129]}\) is \(0.69\).


\section{Summary and Discussion}%
\label{sec:summary}

The different approaches to estimating the distributions in
\Cref{sub:histogram_density_estimate}, \ref{sub:kernel_density_estimate} and \ref{sec:parametric_estimates}
are mutually consistent as can be seen in
\Cref{fig:gumbel_mixture_fits}. This is an assurance of their respective correctness (although they are based one on the
other which slightly undercuts this). The different estimates are also consistent with the distribution-free lower bounds
developed in \Cref{sec:distribution-free_estimates}. This means that selection of the distribution estimate does not
change the objective estimates by a large amount and we can select any one of them.

As stated above, the distribution of \(R_\text{max}\) does not seem to be stationary and this should be further
investigated. Another hint for the non-stationarity can be observed in \Cref{fig:seasonality_of_discharges}, where it
seems the measurements from the beginnings of the years are growing with the pass of time in the mean.

Based on the Q-Q plot in \Cref{fig:gumbel_mixture_fits}, the Gumbel mixture is a good fit. Therefore I select the 
parametric model estimates. This is further
supported by the fact that the Gumbel distribution is often used in flood forecasting models. It is an interesting
conclusion that the river seems to have two modes of functioning. The final estimates are summarized in the table
\Cref{tab:summary_estimates}.

For a regulator that is planning on making an anti-flood barrier on the river near the Niamey facility the recommended
height is dependent on the planned lifetime of this barrier as well as the other possibilities to avoid a disastrous
level of water. Also the price of a meter of the barrier is presumably not constant as a variable of its height. This
should be taken into consideration. I would suggest the 100-year water level (\(\qty{2547.7}{\meter\cubed\per\second}\))
and the barrier height should be derived from this value, but only after further investigation into the stationarity.

\begin{table}[hbt]
    \centering
    \begin{tabularx}{\linewidth}{l>{\centering\arraybackslash}X}
\toprule  Objective & Estimate\\\midrule
  \ref{obj:I} & $\hat{P}[R_{\max} > C] = 0.23$\\\addlinespace
  \ref{obj:II} & $\hat{P}[R_{\max} > D] = 0.0026$\\\addlinespace
  \ref{obj:III} & $\hat{P}[R_{\max} < E] = 0.042$\\\addlinespace
  \ref{obj:IV} - 10-year & $\hat{h}_{0.1} = \qty{2175.9}{\metre\cubed\per\second}$\\\addlinespace
  \ref{obj:IV} - 25-year & $\hat{h}_{0.04} = \qty{2326.2}{\metre\cubed\per\second}$\\\addlinespace
  \ref{obj:IV} - 50-year & $\hat{h}_{0.02} = \qty{2437.4}{\metre\cubed\per\second}$\\\addlinespace
  \ref{obj:IV} - 100-year & $\hat{h}_{0.01} = \qty{2547.7}{\metre\cubed\per\second}$\\\addlinespace
  \ref{obj:IV} - 250-year & $\hat{h}_{0.004} = \qty{2692.8}{\metre\cubed\per\second}$\\\addlinespace
  \ref{obj:V} & $\hat{P}[R_{\max} > F] = 0.005$\\\addlinespace
  \ref{obj:VI} & $\hat{P}[R_{\max} \ge Z + C \mid R_{\max} \ge C] = 0.15$\\\bottomrule
\end{tabularx}
    \caption{Final estimates of the objective values based on the Gumbel mixture parametric estimates.}
    \label{tab:summary_estimates}
\end{table}
